\documentclass[10pt]{article}
\usepackage[margin=0.8in]{geometry}
\usepackage{booktabs,enumitem,graphicx,hyperref,amsmath,url,float,caption}
\captionsetup{font=small,labelfont=bf}
\setlist{itemsep=1.5pt,topsep=3pt}
\hypersetup{colorlinks=true,linkcolor=black,urlcolor=blue!50!black}
\title{\bfseries Cloud-droplet sedimentation and precipitation suppression\\ in the Nd--LWP response of DP-EAMxx (DYCOMS-II RF02)}
\author{Prepared by Claude Code (AI assistant) on behalf of Naser Mahfouz\\ for Naser Mahfouz and Johannes M\"ulmenst\"adt}
\date{Technical note, 6 November 2026}
\begin{document}\sloppy
\maketitle

\begin{abstract}\noindent
We ran DP-EAMxx on the nocturnal DYCOMS-II RF02 drizzling stratocumulus over a 216-run matrix:
grid spacing (40\,km to 1.25\,km) $\times$ droplet number (20--1000\,cm$^{-3}$) $\times$ a cloud-droplet sedimentation multiplier
($\times$0.5, 1, 2) $\times$ precipitation suppression (KK2000 Nd dependence on or off).
The aim is to separate the two microphysical pathways through which Nd controls LWP.
At $\geq$5\,km the model behaves like a single column. There, the positive LWP response to Nd comes entirely from precipitation suppression, and sedimentation acts as a switch
that halves entrainment at low Nd. At 2.5 and 1.25\,km, where boundary-layer cells are resolved, entrainment rises with Nd and falls with sedimentation, smoothly and largely
regardless of the suppression setting. This offsets precipitation suppression, so LWP saturates above Nd$\approx$100--200\,cm$^{-3}$.
At 1.25\,km the response from Nd 200 to 1000 turns slightly negative ($S$ from $-0.003$ to $-0.04$).
None of the runs is in steady state at hour 6, so these are transient (6\,h) results.
\end{abstract}

\section*{Key results}
\begin{enumerate}
\item \textbf{Single-column limit ($\geq$5\,km): suppression explains all of the LWP response.} With suppression off, LWP is flat at 170--175\,g\,m$^{-2}$ ($S\approx0$)
from Nd 50 to 1000 (Figs.~\ref{fig:resp}, \ref{fig:S}). With suppression on, LWP rises monotonically
($S$ = 0.15 for Nd 100$\to$200 and 0.03--0.07 for Nd 200$\to$1000) and has not saturated by Nd 1000 (Fig.~\ref{fig:path}a).
\item \textbf{Sedimentation acts as a threshold at coarse resolution.} Below a cloud-top fall-speed proxy of
sed$\times(r_{\rm eff,top}/10\,\mu{\rm m})^2\approx3.5$ (the switch falls between 3 and 4), $w_e$ is pinned at 4.8--5.2\,mm\,s$^{-1}$. Above it, $w_e$ jumps to 2.5--2.9\,mm\,s$^{-1}$
(Fig.~\ref{fig:we}a), whether suppression is on or off. Low Nd with sed$\times$1 or $\times$2 crosses the threshold, and sed$\times$0.5 never does.
With suppression off this produces a \emph{negative} low-Nd branch ($S_{20\to50}$ from $-0.3$ to $-0.5$), which is the classical
sedimentation--entrainment effect acting as a regime switch.
\item \textbf{Resolved cells ($\leq$2.5\,km): entrainment varies smoothly with sedimentation.} $w_e$ falls almost monotonically with the fall-speed proxy across
all Nd and sed values (Fig.~\ref{fig:we}b), from $\approx$7\,mm\,s$^{-1}$ to $\approx$3\,mm\,s$^{-1}$. At the same proxy, suppression-off runs
entrain 0.3--0.6\,mm\,s$^{-1}$ less.
\item \textbf{At km scale, LWP saturates at high Nd.} The suppression gain levels off at +16 to +23\,g\,m$^{-2}$ beyond Nd 200,
against +45\,g\,m$^{-2}$ and still rising at $\geq$5\,km (Fig.~\ref{fig:path}a). $S_{200\to1000}$ ranges from $-0.04$ to $+0.01$, and at 1.25\,km it is
negative for every sed value. $w_e$ still rises from Nd 200 to 1000 (at 1.25\,km, sed$\times$1: 6.55$\to$6.94\,mm\,s$^{-1}$), so this
looks like the onset of the entrainment-driven falling branch. It is weak at 6\,h.
\item \textbf{The low-Nd branch steepens with sedimentation at km scale.} With suppression on, $S_{20\to50}$ = 0.45/0.57/0.71 (2.5\,km) and 0.53/0.57/0.70 (1.25\,km) for sed $\times$0.5/$\times$1/$\times$2.
Faster sedimentation thins and breaks the Nd-20 deck even with suppression off (cloud cover 0.63--0.69, LWP SD up to 117\,g\,m$^{-2}$; Fig.~\ref{fig:pdf}).
\item \textbf{The runs are transient.} dLWP/dt over hours 4--6 is +10 to +24\,g\,m$^{-2}$\,h$^{-1}$ for most runs at $\geq$5\,km. At 1.25\,km it is about +7\,g\,m$^{-2}$\,h$^{-1}$ at high Nd with suppression on (Fig.~\ref{fig:trend}).
A 12\,h forcing is needed before treating any $S$ value as an equilibrium susceptibility.
\end{enumerate}

\section{Setup}
\paragraph{Case and model.} DYCOMS-II RF02 nocturnal drizzling stratocumulus (6\,h; shortwave off; $T$, $q$ nudged above 650\,hPa;
IOP surface fluxes and large-scale subsidence), run in doubly periodic EAMxx (DP-EAMxx: SHOC, P3, SPA, RRTMGP; 128 levels;
stub sea ice so that one executable runs every resolution). Domains are 600\,km at 40\,km, 300\,km at 20\,km and 240\,km from 10\,km down.
Time steps follow the earlier ACI $\times$ resolution study, and 3D turbulence is off. E3SM: branch \texttt{mahf708/eamxx/p3-cloud-sed-factor}
(E3SM PR \#9017 plus a P3 sedimentation knob).

\paragraph{Matrix.} Every combination of the following, 216 runs in all.
\begin{center}\small
\begin{tabular}{@{}ll@{}}
\toprule
factor & levels \\ \midrule
$\Delta x$ & 40, 20, 10, 5, 2.5, 1.25\,km \\
Nd target (cm$^{-3}$) & 20, 50, 100, 200, 500, 1000 (diagnosed: 10--25\,\% higher) \\
\texttt{p3::cloud\_sedimentation\_factor} & 0.5, 1, 2 \\
precip suppression & on (KK2000, $N_d^{-1.79}$); off ($N_d^{0}$, prefactor $1350\cdot63^{-1.79}=0.812$) \\ \bottomrule
\end{tabular}
\end{center}
The new P3 runtime option \texttt{cloud\_sedimentation\_factor} (default 1) multiplies the mass- and number-weighted cloud-droplet
fall speeds in P3's cloud sedimentation. It is the counterpart of MG2's \texttt{cld\_sed}, which scales the Stokes prefactor.
With precip suppression off, autoconversion no longer depends on Nd, and its rate equals the default rate at Nd = 63\,cm$^{-3}$,
the geometric mean of the original 20--200 sweep. Nd still acts through sedimentation, evaporation and longwave radiation
(shortwave is off).

\paragraph{Diagnostics.} We use the 15-min domain-mean profiles, 15-min per-column 2D fields and 2-hourly 3D snapshots.
AeroCom cloud-top fields come from the new expression operator, \texttt{X.aerocom('top', phase='liq')}.
The LWP spatial variance is computed online as \texttt{((LWP-LWP.mean('col'))$^2$).mean('col')}. Its square root, taken offline, is exact for instantaneous output.
$w_e = dz_i/dt - w_{\rm subs}(z_i)$ is fitted from hour 2 on, with $z_i$ at the $\theta_l$ mid-inversion and $w_{\rm subs}$ from the IOP file.
Cloud cover is the fraction of columns with LWP $>$ 10\,g\,m$^{-2}$.
The susceptibility $S = d\ln{\rm LWP}/d\ln N_d$ uses the diagnosed in-cloud Nd. Unless stated otherwise, values are hour 4--6 means.


\section{Results}
\begin{figure}[H]\centering
\includegraphics[width=\textwidth]{fig1_response.png}
\caption{Hour 4--6 means against diagnosed in-cloud Nd. Columns: resolution. Shade: sedimentation factor. Solid lines with filled circles:
precip suppression on. Dashed lines with open squares: off.}\label{fig:resp}
\end{figure}

\begin{figure}[H]\centering
\includegraphics[width=\textwidth]{fig2_susceptibility.png}
\caption{Susceptibility $S=d\ln{\rm LWP}/d\ln N_d$ (finite differences on diagnosed Nd) on three Nd branches, against resolution.}\label{fig:S}
\end{figure}

\paragraph{Pathway decomposition.} Fig.~\ref{fig:path} isolates each pathway by differencing runs that differ in one switch only.
(a) The suppression pathway, LWP(on) $-$ LWP(off), crosses zero near Nd 63 by construction, since autoconversion rates match there.
It is resolution-independent below Nd 100 and diverges above it: it keeps growing at $\geq$5\,km and saturates at km scale.
(b) The sedimentation pathway, LWP(sed$\times$2) $-$ LWP(sed$\times$0.5), has the opposite sign in the two regimes.
At $\geq$5\,km more sedimentation raises LWP at low Nd (by $\approx$40--55\,g\,m$^{-2}$, through the $w_e$ switch) and lowers it at Nd 100--500 with suppression on.
At km scale it lowers LWP at low Nd (deck break-up) and raises it at Nd 100--500 with suppression on (less entrainment).

\begin{figure}[H]\centering
\includegraphics[width=\textwidth]{figR1_pathways.png}
\caption{(a) Precipitation-suppression pathway at sed$\times$1. (b) Sedimentation pathway with suppression on (solid) and off (dashed). Colour: resolution.}\label{fig:path}
\end{figure}

\paragraph{Sedimentation--entrainment.} P3's cloud-droplet fall speed scales as Stokes, $\propto r^2$, so
sed$\times(r_{\rm eff,top}/10\,\mu$m$)^2$ is a dimensionless proxy for the sedimentation velocity at cloud top. It combines the knob and Nd, through $r_{\rm eff}$.
In Fig.~\ref{fig:we}, every run from both factors lands on one relation per regime. That supports sedimentation removing liquid from the entrainment zone
as the common mechanism behind the sed and Nd effects on $w_e$. The single-column regime responds as a bistable switch, and resolved cells respond smoothly.

\begin{figure}[H]\centering
\includegraphics[width=\textwidth]{figR2_we_vs_vsed.png}
\caption{Entrainment rate against the cloud-top fall-speed proxy, for all 216 runs. (a) $\Delta x\geq$5\,km. (b) $\Delta x\leq$2.5\,km.}\label{fig:we}
\end{figure}

\paragraph{Cloud morphology.} At km scale, low Nd broadens the column-LWP distribution toward thin and clear columns (Fig.~\ref{fig:pdf}).
Faster sedimentation fattens the thin-cloud tail at Nd 20, even with suppression off. At high Nd the distribution narrows and peaks near 200--250\,g\,m$^{-2}$.
LWP spatial SD (Fig.~\ref{fig:resp}) is 1--3\,g\,m$^{-2}$ at $\geq$10\,km, where the columns are nearly identical. It is 6--60\,g\,m$^{-2}$ at 5\,km and 57--117\,g\,m$^{-2}$ at km scale,
largest at low Nd with sed$\times$2.

\begin{figure}[H]\centering
\includegraphics[width=\textwidth]{figR4_lwp_pdf.png}
\caption{Distribution of column LWP over hours 4--6 at 2.5 and 1.25\,km, for Nd 20, 100 and 1000. The area under each curve is the fraction of columns with LWP $>$ 1\,g\,m$^{-2}$.}\label{fig:pdf}
\end{figure}

\paragraph{Transience.} Fig.~\ref{fig:trend} shows the linear LWP trend over the averaging window, and Fig.~\ref{fig:ts} the trajectories.
The coarse runs are still deepening and moistening at hour 6. Only the km-scale runs with suppression off are close to steady.

\begin{figure}[H]\centering
\includegraphics[width=\textwidth]{figR3_lwp_trend.png}
\caption{Linear LWP trend over hours 4--6. The grey band marks $\pm$5\,g\,m$^{-2}$\,h$^{-1}$.}\label{fig:trend}
\end{figure}

\begin{figure}[H]\centering
\includegraphics[width=0.92\textwidth]{fig3_timeseries.png}
\caption{15-min time series at 40\,km and 1.25\,km for Nd 20 and 1000. The shaded band is the hour 4--6 averaging window.}\label{fig:ts}
\end{figure}

\section{Caveats and next steps}
\begin{itemize}
\item \textbf{Transient, single realization.} Run the 12\,h RF02 forcing (\texttt{DYCOMSrf02\_12hr\_PMCnBLH\_4scam.nc}) and 2--3 perturbation seeds
for the km-scale cells, especially Nd 200--1000, where the falling branch is just appearing.
\item \textbf{Sub-km.} 625, 312 and 156\,m are on hold. A targeted set would be Nd \{20, 200, 1000\} $\times$ sed \{0.5, 2\} $\times$ suppression \{on, off\}
at 625\,m (12 runs, $\sim$15 node-h at 6\,h) before committing to 312 and 156\,m.
\item \textbf{Nd targeting.} The SPA factors give diagnosed Nd 10--25\,\% above target. All analysis uses the diagnosed Nd.
\item \textbf{Diagnostics.} Surface precipitation includes cloud-droplet sedimentation. $w_e$ comes from a linear $z_i$ fit from hour 2 on, which is crude for the coarse runs.
A cloud-top sedimentation flux diagnostic (P3 \texttt{qc\_sed} at $z_i$) would replace the proxy in Fig.~\ref{fig:we}.
\end{itemize}

\section{Data and code (Perlmutter, world-readable)}
\begin{itemize}\small
\item Runs: \path{/pscratch/sd/m/mahf708/dp_screamxx_sedps/sedps_<res>_nd<Nd>_sed<0p5|1|2>_ps<on|off>/run/}.
\begin{itemize}
\item Streams: \texttt{horiz\_avg} (profiles), \texttt{aci2d} (columns), \texttt{aerocom} (cloud top and \texttt{lwp\_var}), \texttt{aci3d} (2-hourly 3D).
\end{itemize}
\item Analysis: \path{/pscratch/sd/m/mahf708/dp_screamxx_sedps/analysis/}.
\begin{itemize}
\item \texttt{summary\_h4-6.csv}: one row per run.
\item \texttt{timeseries\_15min.csv}: one row per run and step; run settings as columns.
\item \texttt{susceptibility.csv} and \texttt{lwp\_trend\_h4-6.csv}.
\end{itemize}
\item Scripts: \path{/pscratch/sd/m/mahf708/scmlib/sed-precsupp-study/}.
\begin{itemize}
\item \texttt{analysis/}: Python.
\item \texttt{R/}: R loaders and an example.
\item \texttt{report/}: this report.
\end{itemize}
\item R environment: \path{/pscratch/sd/m/mahf708/r_env/bin/Rscript}. Quick start: \texttt{source(".../R/sedps.R"); runs <- sedps\_runs()}.
\item Code: E3SM \href{https://github.com/mahf708/E3SM/tree/mahf708/eamxx/p3-cloud-sed-factor}{\texttt{mahf708/eamxx/p3-cloud-sed-factor}}
and scmlib \href{https://github.com/mahf708/scmlib/tree/sed-precsupp-study}{\texttt{sed-precsupp-study}}.
\end{itemize}

\end{document}
